\documentclass[10pt,a4paper]{amsart}
\usepackage[margin=19mm]{geometry}
\usepackage{amsmath,amssymb,mathtools}
\usepackage[scr=rsfs,cal=euler]{mathalfa}
\usepackage{xfrac}
\usepackage[unicode,hidelinks,bookmarks=false]{hyperref}
\newcommand{\ie}{i.e.\ }
\newcommand{\cf}{cf.\ }
\newcommand{\resp}{respectively\ }
\newcommand{\Subsection}{\subsection}
\providecommand{\nobreakdash}{\nobreak\hbox{-}}
\setlength{\emergencystretch}{2em}
\multlinegap0pt
\usepackage{amssymb}
\usepackage{enumerate}
\usepackage{cancel}
\usepackage{tikz}
\usepackage{float}
\usepackage{xcolor}
\newtheorem{lemma}{Lemma}
\newtheorem{theorem}{Theorem}
\newcommand{\Real}{\mathbb{R}}
\newcommand{\bR}{{\boldsymbol R}}
\newcommand{\bS}{{\boldsymbol S}}
\newcommand{\ba}{{\boldsymbol a}}
\newcommand{\bb}{{\boldsymbol b}}
\newcommand{\bn}{{\boldsymbol n}}
\newcommand{\bnu}{{\boldsymbol \nu}}
\newcommand{\bt}{{\boldsymbol t}}
\newcommand{\btau}{{\boldsymbol \tau}}
\newcommand{\bv}{{\boldsymbol v}}
\newcommand{\bx}{{\boldsymbol x}}
\newcommand{\by}{{\boldsymbol y}}
\newcommand{\calN}{\mathcal{N}}
\newcommand{\calT}{\mathcal{T}}
\newcommand{\calV}{\mathcal{V}}
\newcommand{\disr}{r_\calT}
\newcommand{\dx}{\,{\rm dx}}
\newcommand{\tr}{^{\mathrm{T}}}
\newcommand{\vertex}{\calV_\calT}
\title{An energy-decreasing algorithm for the finite element approximation of ferronematic equilibrium states}
\author{Alexandre Ern, Ruma R. Maity}
\date{Source: arXiv:2605.11160v1 (CC BY 4.0)}
\begin{document}
\maketitle

\noindent\emph{Source and licence.} An energy-decreasing algorithm for the finite element approximation of ferronematic equilibrium states. Version arXiv:2605.11160v1 (CC BY 4.0), distributed by arXiv under the Creative Commons Attribution 4.0 International licence, http://creativecommons.org/licenses/by/4.0/, as recorded in the arXiv licence metadata for this exact version and shown on the abstract page https://arxiv.org/abs/2605.11160v1. Full licence notice: \texttt{licenses/arxiv-2605.11160-CC-BY-4.0.txt}.\par\smallskip

\noindent\textit{Editorial scope.} Counted unit: the article's theorem thm:engy-decr-disc (source0.tex lines 498--501). The article's complete original proof of it is delivered with it (source0.tex lines 503--537, one \texttt{proof} environment of 35 lines). No mathematics is written, repaired or re-derived: the statement and the proof are the article's own lines, in the article's own notation, and no retained line is substituted or re-flowed.  Retained uncounted as the source-local context the proof needs: lines 291--294; lines 361--363; lines 391--398; lines 419--421; lines 423--437; lines 459--461; lines 474--479.  Nothing was omitted from any retained span, so the package carries no omission record.  No retained line contains a citation, so the wrapper carries no bibliography and no bibliography entry.  No retained line contains a figure environment, an image-inclusion command or drawing code, so no image is delivered and the package declares no asset dependency.  Editorial formatting only, in addition: an AMS \texttt{amsart} preamble that restates the article's own declarations and macros used by the retained lines, namely the environments \texttt{lemma}, \texttt{theorem} on separate counters, together with the standard packages those lines need.  The retained environments print in this document as Lemma 1, Theorem 1 in order of appearance, whereas the article prints the same items with the numbering of its own full document.  The complete native source of the article as distributed in the arXiv e-print for this exact version is bundled unchanged as \texttt{sources/200343/source0.tex}.
\begin{lemma}[Projection onto sphere] \label{lem:PS1}
\textup{(i)} $\|P_{\bS}(\bx)-P_{\bS}(\by)\|_{\Real^2} \le \|\bx-\by\|_{\Real^2}$ for all $\bx,\by \in \Real^2_{\ge1}$.
\textup{(ii)} For all $\bv \in H^1(\Omega;\Real^2)$ taking values in $\Real^2_{\ge1}$, $P_{\bS}(\bv) \in H^1(\Omega;\Real^2)$ and takes values in $\bS$; moreover, $\|\nabla P_{\bS}(\bv)\|_{\Real^{2\times2}} \le \|\nabla\bv\|_{\Real^{2\times2}}$. 
\end{lemma}

\begin{equation} \label{eq:DMP}
k_{\ba\bb} \le 0 \quad \forall \ba,\bb \in \vertex, \; \ba\ne\bb.
\end{equation}

\begin{lemma}[Energy evaluation] \label{lem:engy_eval}
The following holds for all $\Phi_\calT \in W_\calT$:
\begin{equation} \label{eq:engy_eval}
\int_\Omega \|\nabla \Phi_\calT\|_{\Real^{4\times2}}^2 \dx = - \frac12
\sum_{\ba,\bb\in\vertex} k_{\ba\bb}\Big\{ \, Q_c^2\|\bnu(\ba)-\bnu(\bb)\|_{\Real^2}^2
+ M_c^2\|\bn(\ba)-\bn(\bb)\|_{\Real^2}^2 \, \Big\}.
\end{equation}
\end{lemma}

\begin{equation} \label{eq:def_phi}
\varphi(r) := \frac{2r}{1-r^2} \quad \forall r\in(-1,1).
\end{equation}

\begin{lemma}[Geometry of $\calN$] \label{lem:phi}
Let $\bn\in\bS$ and let $\bt\in T_{\bn}\bS$ be the unit vector 
obtained from $\bn$ by a rotation
of $\frac{\pi}{2}$ in $\Real^2$. Set $\bnu:=\bR(\bn)\in\bS$ and let 
$\btau\in T_{\bnu}\bS$ be the unit vector obtained from $\bnu$ by a rotation
of $\frac{\pi}{2}$ in $\Real^2$. The following holds for all $r\in (-1,1)$:
\begin{equation} \label{eq:RPS}
(\bR \circ P_{\bS})(\bn + r \bt) = P_{\bS}(\bnu + \varphi(r)\btau),
\end{equation}
with the function $\varphi:(-1,1)\rightarrow\Real$ defined in~\eqref{eq:def_phi}. 
In other words, we have
\begin{equation} \label{eq:RPinN}
\big(Q_c P_{\bS}(\bnu + \varphi(r)\btau),M_cP_{\bS}(\bn + r \bt)\big)\tr \in \calN.
\end{equation}
\end{lemma}

\begin{equation} \label{eq:enrg_decr_calT}
E(\Psi_\calT^{j+1}) \le E(\Psi_\calT^{j}).
\end{equation}

\begin{align}
F^j(\disr) := -\frac12 \sum_{\ba,\bb\in\vertex} k_{\ba\bb} \Big\{ \,
&Q_c^2\|(\bnu^j(\ba)+\varphi(r_{\ba})\btau^j(\ba))-(\bnu^j(\bb)+\varphi(r_{\bb})\btau^j(\bb))\|_{\Real^2}^2 \nonumber \\
+ {} &M_c^2\|(\bn^j(\ba)+r_{\ba}\bt^j(\ba))-(\bn^j(\bb)+r_{\bb}\bt^j(\bb))\|_{\Real^2}^2
\, \Big\}. \label{eq:def_Fj}
\end{align} 

\begin{theorem}[Energy decay] \label{thm:engy-decr-disc}
Assume that the triangulation is strictly acute so that \eqref{eq:DMP} holds true.
Then the sequence $(\Psi_\calT^j)_{j=1}^\infty$ satisfies~\eqref{eq:enrg_decr_calT}. 
\end{theorem}

\begin{proof}
Owing to Lemma~\ref{lem:phi}, we infer that
\[
\bnu^{j+1}(\ba):= \bR(\bn^{j+1}(\ba)) = (\bR \circ P_{\bS})(\bn^j(\ba)+r^*_{\ba}\bt^j(\ba))
=P_{\bS}(\bnu^j(\ba)+\varphi(r^*_{\ba})\btau^j(\ba)).
\]
Invoking Lemma~\ref{lem:engy_eval} gives
\begin{align*}
\int_\Omega &\|\nabla \Psi_\calT^{j+1}\|_{\Real^{4\times2}}^2 \dx \\= {}& - \frac12
\sum_{\ba,\bb\in\vertex} k_{\ba\bb}\Big\{ \, Q_c^2\|\bnu^{j+1}(\ba)-\bnu^{j+1}(\bb)\|_{\Real^2}^2
+ M_c^2\|\bn^{j+1}(\ba)-\bn^{j+1}(\bb)\|_{\Real^2}^2 \, \Big\} \\
={}&  - \frac12
\sum_{\ba,\bb\in\vertex} k_{\ba\bb}\Big\{ \, Q_c^2\|P_{\bS}(\bnu^j(\ba)+\varphi(r^*_{\ba})\btau^j(\ba)) - P_{\bS}(\bnu^j(\bb)+\varphi(r^*_{\bb})\btau^j(\bb))\|_{\Real^2}^2 \\
&+ M_c^2 \|P_{\bS}(\bn^j(\ba)+r^*_{\ba}\bt^j(\ba))-P_{\bS}(\bn^j(\bb)+r^*_{\bb}\bt^j(\bb))\|_{\Real^2}^2 \, \Big\}.
\end{align*}
Invoking Lemma~\ref{lem:PS1} and crucially that $k_{\ba\bb}\le 0$ (see~\eqref{eq:DMP}),
we infer that
\begin{align*}
\int_\Omega &\|\nabla \Psi_\calT^{j+1}\|_{\Real^{4\times2}}^2 \dx \\
\le{}&  - \frac12
\sum_{\ba,\bb\in\vertex} k_{\ba\bb}\Big\{ \, Q_c^2\|(\bnu^j(\ba)+\varphi(r^*_{\ba})\btau^j(\ba)) - (\bnu^j(\bb)+\varphi(r^*_{\bb})\btau^j(\bb))\|_{\Real^2}^2 \\
&+ M_c^2 \|(\bn^j(\ba)+r^*_{\ba}\bt^j(\ba))-(\bn^j(\bb)+r^*_{\bb}\bt^j(\bb))\|_{\Real^2}^2 \, \Big\}
\end{align*}
Recalling~\eqref{eq:def_Fj}, this proves that
\[
\int_\Omega \|\nabla \Psi_\calT^{j+1}\|_{\Real^{4\times2}}^2 \dx \le F^j(\disr^*).
\]
Observing that $\int_\Omega \|\nabla \Psi_\calT^{j}\|_{\Real^{4\times2}}^2 \dx = F^j(0)$,
we conclude that
\[
\int_\Omega \|\nabla \Psi_\calT^{j+1}\|_{\Real^{4\times2}}^2 \dx \le F^j(\disr^*)
\le F^j(0) = \int_\Omega \|\nabla \Psi_\calT^{j}\|_{\Real^{4\times2}}^2.
\]
This proves the claim.
\end{proof}

\end{document}
